\documentclass[../../main.tex]{subfiles}% 注意这里的文件路径不能用 ./main.tex ,否则用latexmk编译子文件会报错
\graphicspath{{\subfix{./image/}}{\subfix{../../image/}}} % 指定图片引用路径,后续可以直接使用图片文件名
% 如果图片存放在上级目录,则使用 ../../image/ 作为备用路径

% 例如:
% \begin{figure}[H]
% \centering
% \includegraphics[width=0.3\textwidth]{图.png}
% \caption{}
% \label{figure:图}
% \end{figure}
% 注意:上述\label{}一定要放在\caption{}之后,否则引用图片序号会只会显示??.

\begin{document}

\section{能量模估计}

\begin{theorem}
假设 $c(x)\geqslant c_0>0$，$u\in C^2(\Omega)\cap C^1(\overline{\Omega})$ 是 $n$ 维位势方程的 Dirichlet 问题\eqref{eq:dirichlet-energy}
\begin{align}
\begin{cases}
-\Delta u + c(x)u = f(x), & x\in\Omega, \\
u|_{\partial\Omega} = 0.
\end{cases} \label{eq:dirichlet-energy}
\end{align}
的解，则
\begin{align}
\int_\Omega |Du(x)|^2\,\mathrm{d}x + \frac{c_0}{2}\int_\Omega |u(x)|^2\,\mathrm{d}x \leqslant M\int_\Omega |f(x)|^2\,\mathrm{d}x, \label{eq:energy-est-dir}
\end{align}
其中 $M$ 是仅依赖于 $c_0$ 的常数。
\end{theorem}

\begin{proof}
在 Dirichlet 问题 \eqref{eq:dirichlet-energy} 的方程两端同乘 $u$，再在 $\Omega$ 上积分得到
\begin{align*}
-\int_\Omega u\Delta u\,\mathrm{d}x + \int_\Omega c(x)u^2\,\mathrm{d}x = \int_\Omega fu\,\mathrm{d}x.
\end{align*}
对于上式左端第一项应用 \hyperref[Basis of Analytics-theorem:Gauss-Green公式(散度定理)]{Gauss-Green公式}，右端应用 Cauchy 不等式
\begin{align*}
2ab \leqslant \varepsilon a^2 + \frac{1}{\varepsilon}b^2,\quad \varepsilon>0,
\end{align*}
则
\begin{align*}
\int_\Omega |Du|^2\,\mathrm{d}x + c_0\int_\Omega u^2\,\mathrm{d}x \leqslant \frac{c_0}{2}\int_\Omega u^2\,\mathrm{d}x + \frac{1}{2c_0}\int_\Omega f^2\,\mathrm{d}x.
\end{align*}
上式移项即得证。\end{proof}

\begin{lemma}[Friedrichs 不等式]
假设 $u\in C_0^1(\Omega)$，则
\begin{align}
\int_\Omega |u(x)|^2\,\mathrm{d}x \leqslant 4d^2\int_\Omega |Du(x)|^2\,\mathrm{d}x, \label{eq:friedrichs}
\end{align}
其中 $d$ 是 $\Omega$ 的直径。
\end{lemma}

\begin{proof}
由于 $\Omega$ 的直径为 $d$，则可以做一个边长为 $2d$ 且平行于坐标轴的 $n$ 维正方体将 $\Omega$ 包含于其中。不妨设此正方体为
\begin{align*}
Q = \{x=(x_1,x_2,\cdots,x_n) \mid 0\leqslant x_i\leqslant 2d,\ i=1,2,\cdots,n\}.
\end{align*}
令
\begin{align*}
\tilde{u} = \begin{cases}
u, & x\in\Omega, \\
0, & x\in Q\setminus\Omega,
\end{cases}
\end{align*}
则显然 $\tilde{u}\in C_0^1(Q)$ 且
\begin{align*}
\tilde{u}(x_1,x_2,\cdots,x_n) = \int_0^{x_1} \tilde{u}_\xi(\xi,x_2,\cdots,x_n)\,\mathrm{d}\xi.
\end{align*}
利用 Schwarz 不等式得
\begin{align*}
\tilde{u}^2(x)\leqslant x_1\int_0^{x_1} |\tilde{u}_\xi(\xi,x_2,\cdots,x_n)|^2\,\mathrm{d}\xi \leqslant 2d\int_0^{2d} |\tilde{u}_{x_1}(x_1,x_2,\cdots,x_n)|^2\,\mathrm{d}\xi.
\end{align*}
两端对 $x$ 在 $Q$ 上积分，则
\begin{align*}
\int_Q \tilde{u}^2(x)\,\mathrm{d}x \leqslant 4d^2\int_Q |\tilde{u}_{x_1}|^2\,\mathrm{d}x \leqslant 4d^2\int_Q |D\tilde{u}|^2\,\mathrm{d}x.
\end{align*}
注意到 $\tilde{u}$ 的定义，我们得到不等式 \eqref{eq:friedrichs}。\end{proof}

\begin{theorem}
假设 $c(x)\geqslant 0$。如果 $u\in C^2(\Omega)\cap C^1(\overline{\Omega})$ 是问题 \eqref{eq:dirichlet-energy} 的解，则
\begin{align}
\int_\Omega |Du(x)|^2\,\mathrm{d}x + \int_\Omega |u(x)|^2\,\mathrm{d}x \leqslant M\int_\Omega |f(x)|^2\,\mathrm{d}x, \label{eq:energy-est-dir2}
\end{align}
其中 $M$ 是仅依赖于 $\Omega$ 的直径的常数。
\end{theorem}

下面考虑位势方程的第三边值问题
\begin{align}
\begin{cases}
-\Delta u + c(x)u = f(x), & x\in\Omega, \\
\left[\frac{\partial u}{\partial n} + \alpha(x)u\right]\bigg|_{\partial\Omega} = 0.
\end{cases} \label{eq:third-energy}
\end{align}
同样利用 \hyperref[Basis of Analytics-theorem:Gauss-Green公式(散度定理)]{Gauss-Green公式}我们可以得到下面的能量模估计。

\begin{theorem}
假设 $c(x)\geqslant c_0>0$，$\alpha(x)\geqslant 0$。如果 $u\in C^2(\Omega)\cap C^1(\overline{\Omega})$ 是第三边值问题 \eqref{eq:third-energy} 的解，则
\begin{align}
\int_\Omega |Du(x)|^2\,\mathrm{d}x + \frac{c_0}{2}\int_\Omega |u(x)|^2\,\mathrm{d}x + \int_{\partial\Omega} \alpha(x)u^2(x)\,\mathrm{d}S(x) \leqslant M\int_\Omega |f(x)|^2\,\mathrm{d}x, \label{eq:energy-est-third}
\end{align}
其中 $M$ 是仅依赖于 $c_0$ 的常数。
\end{theorem}





\end{document}